\scanpage{13}

\deleted{irréductible en vertu du th. de Matsusaka.} \deleted{\unsure{Alors}} \uline{On suppose déjà} que $A_{\unsure{a}}$ est une jacobienne \inserted{polarisée} ; \deleted{\unsure{alors qu'il existe|donc qu'il est} Or un résultat d'\uline{unicité} de Matsusaka affirme qu'il \unsure{n'}existe alors \inserted{à translation près} qu'un seul div. $D$ positif sur $A_a$ \deleted{\unsure{ayant les}}, \uline{irréductible ou non}, qui \deleted{ait les} satisfasse à deux équivalences numériques comme ci-dessus, et qu'alors} ce diviseur $\Theta_a$ est dans la classe de polarisation donc est \uline{irréductible} ; \deleted{donc} $C_a$ est irréductible et non \unsure{singulière|singulier} sur $k$. Donc $C$ se \deleted{réduit bien, et il} est simple sur $S$, et on achève maintenant \unsure{trivialement|facilement}.

La démonstration donne d'ailleurs ceci :

\deleted{Soit} \uline{Corollaire} 1 Soit $T$ \deleted{l'image} $=\overline{f(M)}$, où $f$ est le morphisme structural $M\to S$. Pour que $t\in T$ soit ds $f(M)$, il f. et s. que le diviseur de polarisation de $A_t$ [défini à translation près, en vertu de $\deg\Theta^g/g!=1$] soit \inserted{\unsure{ablt.|relt.|vraiment}} irréductible. \deleted{\unsure{§}}

% archive sheet number: 12
